\section{Інженерний процес TDD, критерії прийняття та BDD-сценарії}\label{sec:tdd}

Цей розділ фіксує обов'язковий інженерний регламент розробки програмного забезпечення за методологією Test-Driven Development (TDD) та формальні верифікаційні критерії прийняття.

\subsection{Обов'язковий життєвий цикл розробки TDD}\label{sec:tdd-process}

\begin{figure}[htbp]
\centering
\begin{tikzpicture}[node distance=8mm and 10mm]
  \node[sked box] (t1) {1. Написання\\тесту};
  \node[sked box, right=of t1] (t2) {2. Перевірка\\падіння (Red)};
  \node[sked box, right=of t2] (t3) {3. Реалізація\\коду (Green)};
  \node[sked box, below=of t3] (t4) {4. Рефакторинг\\коду};
  \node[sked box, left=of t4] (t5) {5. Повторний\\запуск тестів};
  
  \draw[sked arrow] (t1) -- (t2);
  \draw[sked arrow] (t2) -- (t3);
  \draw[sked arrow] (t3) -- (t4);
  \draw[sked arrow] (t4) -- (t5);
  \draw[sked back] (t5) to[out=180,in=180,looseness=1.5] node[sked label, left=2pt] {Наступна функція} (t1);
\end{tikzpicture}
\caption{Обов'язковий 5-кроковий цикл розробки за TDD}\label{fig:tdd-cycle}
\end{figure}

\begin{itemize}[leftmargin=*,labelindent=0pt]
    \item \textbf{REQ-TDD-001 (Обов'язковість процесу TDD):} Розробка будь-якої нової функціональності, бізнес-правил та виправлення дефектів повинні здійснюватися суворо за методологією TDD.~\sked{K7,K11}
    \textit{Умова верифікації:} Будь-який комміт або пул-реквест у репозиторій повинен містити відповідні автоматизовані тести, створені суворо перед написанням коду реалізації.

    \item \textbf{REQ-TDD-002 (5-кроковий цикл розробки):} Кожен етап розробки виконується за чіткою послідовністю кроків: (1) створити автоматизований тест для очікуваної поведінки перед написанням функціонального коду; (2) переконатися, що тест падає саме з очікуваної причини (відсутність функції/неправильний результат); (3) написати мінімальний код для проходження тесту; (4) провести рефакторинг коду; (5) повторно запустити відповідний набір тестів і переконатися в їх успішності.~\sked{K12}
    \textit{Умова верифікації:} Історія змін та логи запусків тестів підтверджують проходження стадій «Red $\rightarrow$ Green $\rightarrow$ Refactor».
\end{itemize}

\subsection{Сфера тестування та перевірка граничних випадків}\label{sec:tdd-scope}

\begin{itemize}[leftmargin=*,labelindent=0pt]
    \item \textbf{REQ-TDD-003 (Покриття бізнес-правил та граничних випадків):} Автоматизовані тести повинні верифікувати всі погоджені в цій SRS бізнес-правила, обмеження ролей та суттєві граничні випадки системи бронювання.~\sked{K13}
    \textit{Обов'язкові граничні тест-кейси:}
    \begin{itemize}
        \item Бронювання рівно 1 місця (нижня межа) та рівно 6 місць (верхня межа);
        \item Спроба бронювання 0 та 7 місць (порушення меж діапазону);
        \item Конкурентна спроба забронювання одного й того самого місця кількома паралельними потоками;
        \item Запит на бронювання кількох місць, коли одне з них уже зайняте;
        \item Скасування за 1 секунду до часу початку сеансу (успіх) та в момент початку (відхилення);
        \item Повторний виклик скасування вже скасованого бронювання (перевірка ідемпотентності);
        \item Спроба модифікації сеансу з активними бронюваннями (блокування) та скасування сеансу без введення причини (блокування).
    \end{itemize}
    
    \item \textbf{REQ-TDD-004 (Відсутність жорсткого порогу покриття):} Конкретний фіксований відсоток покриття рядків коду тестами (coverage threshold) на поточному етапі не встановлюється; пріоритетом є якісна повнота перевірки бізнес-сценаріїв та граничних умов.~\sked{K14}
    
    \item \textbf{REQ-TDD-005 (Критерії до тестового інструментарію):} Фреймворк тестування обирається разом із технологічним стеком, проте обов'язково повинен підтримувати швидкий запуск юніт- та інтеграційних тестів, можливість емуляції конкурентних запитів та генерацію структурованих звітів запусків.~\sked{K15,K43}
\end{itemize}

\subsection{Критерії прийняття та BDD-сценарії поведінки системи}\label{sec:tdd-bdd}
Для ключових бізнес-правил формулюються обов'язкові сценарії мовою Gherkin (Given-When-Then), що слугують безпосередньою основою для написання тестів~\sked{K13}.

\subsubsection{Сценарій 1: Конкурентне бронювання одного місця}\label{sec:bdd-sc1}
\begin{skedcode}
Сценарій: Конкурентне бронювання одного й того самого місця двома користувачами
  Дано: Користувач_1 та Користувач_2 автентифіковані у системі
    Та місце "Ряд 3, Місце 5" на сеанс ID=101 є вільним
  Коли: Користувач_1 надсилає запит на бронювання місця "Ряд 3, Місце 5"
    Та одночасно Користувач_2 надсилає запит на бронювання того ж місця
  Тоді: Рівно одне бронювання створюється зі статусом "Активне"
    Та інший користувач отримує помилку про зайнятість місця
    Та місце залишається зайнятим і не бронюється повторно
\end{skedcode}
\noindent\textit{Пов'язані вимоги:} REQ-F-008, REQ-F-009, REQ-NF-001~\sked{K29}.

\subsubsection{Сценарій 2: Відхилення бронювання при зайнятості хоча б одного місця}\label{sec:bdd-sc2}
\begin{skedcode}
Сценарій: Бронювання трьох місць, коли одне з них зайняте (все або нічого)
  Дано: Користувач автентифікований у системі
    Та на сеанс ID=101 місця 1 і 2 вільні, а місце 3 уже заброньоване іншим клієнтом
  Коли: Користувач надсилає запит на бронювання групи місць [1, 2, 3]
  Тоді: Запит відхиляється повністю
    Та жодного запису бронювання для користувача не створюється
    Та місця 1 і 2 залишаються вільними для вибору
\end{skedcode}
\noindent\textit{Пов'язані вимоги:} REQ-F-010, REQ-NF-001~\sked{K30}.

\subsubsection{Сценарій 3: Успішне скасування бронювання до початку сеансу}\label{sec:bdd-sc3}
\begin{skedcode}
Сценарій: Успішне скасування бронювання до початку сеансу
  Дано: Користувач автентифікований у системі
    Та користувач має початкове активне бронювання ID=501 на сеанс із часом початку 19:00
    Та поточний системний час становить 18:45 (менший за час початку сеансу)
  Коли: Користувач ініціює скасування власного бронювання ID=501
  Тоді: Бронювання переходить у статус "Скасоване"
    Та всі зарезервовані за ним місця негайно стають вільними
\end{skedcode}
\noindent\textit{Пов'язані вимоги:} REQ-F-011, REQ-F-012, REQ-F-013~\sked{K36,K37,K38}.

\subsubsection{Сценарій 4: Блокування скасування бронювання точно в момент початку сеансу}\label{sec:bdd-sc4}
\begin{skedcode}
Сценарій: Блокування скасування точно в момент настання часу початку сеансу
  Дано: Користувач автентифікований у системі
    Та користувач має початкове активне бронювання ID=502 на сеанс із часом початку 19:00:00
    Та поточний системний час становить рівно 19:00:00 (настав час початку сеансу)
  Коли: Користувач надсилає запит на скасування власного бронювання ID=502
  Тоді: Система блокує операцію з повідомленням про початок сеансу
    Та статус бронювання ID=502 залишається "Активне"
    Та зарезервовані місця не звільняються
\end{skedcode}
\noindent\textit{Пов'язані вимоги:} REQ-F-012~\sked{K37}.

\subsubsection{Сценарій 5: Блокування скасування бронювання після початку сеансу}\label{sec:bdd-sc5}
\begin{skedcode}
Сценарій: Блокування спроби скасування бронювання після початку сеансу
  Дано: Користувач автентифікований у системі
    Та користувач має початкове активне бронювання ID=503 на сеанс із часом початку 19:00
    Та поточний системний час становить 19:15 (сеанс уже триває)
  Коли: Користувач намагається скасувати власне бронювання ID=503
  Тоді: Система блокує операцію з повідомленням про те, що сеанс уже триває
    Та статус бронювання ID=503 залишається "Активне"
    Та зарезервовані місця не звільняються
\end{skedcode}
\noindent\textit{Пов'язані вимоги:} REQ-F-012~\sked{K37}.

\subsubsection{Сценарій 6: Блокування зміни та явне скасування сеансу адміністратором}\label{sec:bdd-sc6}
\begin{skedcode}
Сценарій: Обробка сеансу з активними бронюваннями адміністратором
  Дано: Адміністратор увійшов до адміністративного розділу
    Та на сеанс ID=202 існують початкові активні бронювання користувачів
  Коли: Адміністратор намагається змінити час або зал сеансу ID=202
  Тоді: Система блокує зміну параметрів сеансу через наявність активних бронювань
  Коли: Адміністратор обирає "Скасувати сеанс", вводить причину "Технічна несправність проєктора" і підтверджує
  Тоді: Сеанс ID=202 отримує статус "Скасовано"
    Та всі активні бронювання на сеанс ID=202 переходять у статус "Скасовано"
    Та в кабінеті кожного клієнта відображається причина "Технічна несправність проєктора"
\end{skedcode}
\noindent\textit{Пов'язані вимоги:} REQ-F-017, REQ-F-018, REQ-F-019, REQ-F-020~\sked{K45,K46}.
